\begin{tikzpicture}[scale=3]
    % Definiranje koordinat oglišč kocke
    % Sprednja ploskev (v_point - prazni krogci)
    \coordinate (A) at (0,0,0);
    \coordinate (B) at (1,0,0);
    \coordinate (C) at (1,1,0);
    \coordinate (D) at (0,1,0);
    
    % Zadnja ploskev (vnode - polne pike)
    % Zamaknjena v x, y in z smeri za 3D učinek
    \def\zshift{0.4}
    \coordinate (A') at ($(A)+(\zshift,-\zshift)$);
    \coordinate (B') at ($(B)+(\zshift,-\zshift)$);
    \coordinate (C') at ($(C)+(\zshift,-\zshift)$);
    \coordinate (D') at ($(D)+(\zshift,-\zshift)$);

    % ---------------- RISANJE ROBOV ----------------
    
    % Zadnja ploskev (črtkano za 3D učinek, ali polno, če želite strogo samo new_edge)
    % Tukaj uporabljam dashed_edge kot variacijo new_edge za preglednost.
    % Če želite nujno samo "new_edge", zamenjajte dashed_edge z new_edge.
    \draw[new_edge] (A') -- (B') -- (C') -- (D') -- cycle;

    % Povezave med sprednjo in zadnjo ploskvijo (new_edge)
    \draw[new_edge] (A) -- (A');
    \draw[new_edge] (B) -- (B');
    \draw[new_edge] (C) -- (C');
    \draw[new_edge] (D) -- (D');

    % Sprednja ploskev (new_edge - polne črte)
    \draw[new_edge] (A) -- (B) -- (C) -- (D) -- cycle;

    % ---------------- RISANJE TOČK ----------------
    
    % Zadnja ploskev - vnode (polne pike)
    \node[v_point] at (A') {};
    \node[v_point] at (B') {};
    \node[v_point] at (C') {};
    \node[v_point] at (D') {};

    % Sprednja ploskev - v_point (prazni krogci)
    \node[v_point] at (A) {};
    \node[v_point] at (B) {};
    \node[v_point] at (C) {};
    \node[v_point] at (D) {};

\end{tikzpicture}